\scanpage{22}
\mspage{- 3 -}
\noindent\hfill n\textsuperscript{o}355

Le $B$-module $E\handwritten{\otimes}_A F$ est réunion des sous-modules $E'\handwritten{\otimes}_A F$, $E'$ parcourant l'ensemble des sous-modules de type fini de $E$. Si donc $\underline P$ appartient à $\operatorname{Ass}_B(E\handwritten{\otimes}_A F)$, il existe un $E'$ de type fini tel que $\underline P\handwritten{\in}\operatorname{Ass}_B(E'\handwritten{\otimes}_A F)$. D'aprè\unsure{s} (ii), il existe $\underline p\handwritten{\in}\operatorname{Ass}_A(E')$ tel que $\underline P\handwritten{\in}\operatorname{Ass}_B(F\handwritten{/}\underline p F)$ ; comme $\operatorname{Ass}_A(E')\handwritten{\subset}\operatorname{Ass}_A(E)$, cela achève de démontrer la proposition.

COROLLAIRE.- \uline{Si} $A$ \uline{est noethérien}, \uline{et si} $\underline P\handwritten{\in}\operatorname{Ass}_B(E\handwritten{\otimes}_A F)$, \uline{on a} $\handwritten{\pi}^{-1}(\underline P)\handwritten{\in}\operatorname{Ass}_A(E)$.

Cela résulte des propositions 1 et 2.

\uline{Remarques}.-

1) Supposons $A$ \uline{et} $B$ \uline{noethériens}. Soit $\underline P$ un idéal premier de $B$, et soit $\underline p=\handwritten{\pi}^{-1}(\underline P)$. Soit $S=A-\underline p$, et soit $k(\underline p)=S^{-1}(A\handwritten{/}\underline p)$ le corps des fractions de $A\handwritten{/}\underline p$. Puisque $\underline P$ contient $\handwritten{\pi}(\underline p)$, il définit un idéal premier de $B\handwritten{/}\underline p B$, et puisqu'il ne rencontre pas $S$, il définit aussi un idéal premier, noté $\underline P'$, de $S^{-1}(B\handwritten{/}\underline p B)=k(\underline p)\handwritten{\otimes}_A B$. Alors, \uline{pour que} $\underline P\handwritten{\in}\operatorname{Ass}_B(E\handwritten{\otimes}_A F)$, \uline{il faut et il suffit que} $\underline p\handwritten{\in}\operatorname{Ass}_A(E)$ \uline{et} $\underline P'\handwritten{\in}\operatorname{Ass}_{k(\underline p)\handwritten{\otimes}_A B}(k(\underline p)\handwritten{\otimes}_A F)$.

Vu la prop.2, tout revient à voir que les conditions

\[
\text{`` }\underline P\handwritten{\in}\operatorname{Ass}_B(F\handwritten{/}\underline p F)\text{ '' \quad et \quad `` }\underline P'\handwritten{\in}\operatorname{Ass}_{k(\underline p)\handwritten{\otimes}_A B}(k(\underline p)\handwritten{\otimes}_A F)\text{ ''}
\]

sont équivalentes, ce qui résulte de la prop.5 du \handwritten{\S} 1 (c'est ici que l'on utilise l'hypothèse noethérienne sur $B$).

\handwritten{[}D'après Grothendieck, l'énoncé ci-dessus est celui que l'on utilise dans les schémas (le produit tensoriel par $k(\underline p)$ s'interprétant comme le passage à la ``fibre au-dessus de $\underline p$'').\handwritten{]}

2) La proposition 2 s'applique notamment lorsque $B$ \uline{est plat sur} $A$ et que l'on prend $F=B$. Le module $E\handwritten{\otimes}_A F$ est alors simplement le module $E_{(B)}$ \uline{déduit de} $E$ \uline{par extension des scalaires}.

On recopie ensuite les deux dernières pages de l'Etat 1.

\begin{center}-:-:-:-\end{center}

% archive sheet number: 21
